\section{合卷：皇帝退位之后，帝国的问题仍在}

1911至1912年的转折改变了最高政治名义，却不是所有社会关系同时翻页。咨议局、商会、报刊、新军、铁路公司和旧式地方网络既提供动员的渠道，也使各省对税收、债务、兵权和代表资格有不同期待。武昌之后的独立、议和和建国因此是一组相互牵动的地方过程，不是一条从起义直达共和国的直线。

退位诏书解决了谁以皇帝名义统治的制度问题；它没有决定军队如何安置、盐税和外债如何承继、旗人生计如何维持，或县镇如何恢复日常秩序。革命的意义正同时在于它打开了新的政治合法性，也暴露了旧帝国留下的财政、军事和社会基础仍须重新协商。

\subsection{制度得失}

\textbf{所得：} 皇帝与部院的旧中心退出，省级公共政治、共和国名义和新的组织形式获得空间。

\textbf{代价：} 财政、军队和地方秩序缺乏一项能立即替代帝国协调机制的安排，政治参与的扩大仍受资源与地域限制。

\textbf{历史余波：} 清朝的终结不是本卷问题的终结，而是征服、财政、边疆与改革如何在新政体中继续发生争议的起点。
